\section{Optimised codebooks}
\label{app:codebooks}
\subsection*{ConvAbuse -- blind}
\begin{quote}\small
---



Instruction Set: Judging Annotation Uncertainty (Codebook)



## Objective

Your task is to estimate the degree of uncertainty a human annotator would have when labeling the provided item. Specifically, predict the likelihood that an annotator would change their label (between '0' - Non-Abusive, or '1' - Abusive) if they evaluated the same item under different criteria or were evaluating it repeatedly.



**Output Requirements:**

Use the following scale for your estimate:

- **Low Uncertainty:** The annotator would likely assign the same label (0 or 1) consistently.

- **High Uncertainty:** The annotator would likely vary their label assignment across evaluations.



## Dimensions of Uncertainty

To make this judgment, analyze the item based on three dimensions. High Uncertainty is indicated when factors from *any* of these dimensions are present. Low Uncertainty is indicated when factors in *none* of these dimensions are present.



### 1. Clarity of Intent (Literal vs. Subtext)

*   **Definition:** Whether the abusive or safe nature of the text is immediately obvious from the words themselves, or if it relies on subtle tone.

*   **Indicators of Low Uncertainty (High Certainty):**

    *   Clear policy triggers (e.g., explicit slurs, explicit threats, explicit insults).

    *   Clear benign text (e.g., "Hello", "Ok", "I agree", factual questions).

    *   Direct, unambiguous phrasing that avoids irony or metaphors.

*   **Indicators of High Uncertainty:**

    *   **Ambiguous Intent:** Phrasing that could reasonably be interpreted as hostile *or* benign depending on context (e.g., "you're having a problem?", "Why did you do that?").

    *   **Irony/Sarcasm:** Text where the literal meaning differs from the intended tone, making the abuse classification subjective.

    *   **Support for Denial:** Text that appears innocent but might imply underlying harassment or intellectual denigration when read critically.



### 2. Contextual Dependence

*   **Definition:** How much previous conversation history is required to determine the meaning or tone of the current turn.

*   **Indicators of Low Uncertainty:**

    *   **Standalone Meaning:** The current message is understandable and classified without needing to reference previous turns (e.g., "NO", "of course.", "Ok").

    *   **Rigid Context:** Even with history, the current message's intent remains clear and consistent (e.g., if history involves abuse, and current text says "That was mean", the intent is abusive).

*   **Indicators of High Uncertainty:**

    *   **Heavy History Reliance:** Classification fundamentally depends on interpreting the emotional state established in `prev\_agent` / `prev\_user` segments.

    *   **Exogenous Context:** The meaning relies on conversational flow where the current message functions as a tone-setter rather than a new topic.

    *   *Note:* Example 1 shows that text relying on a specific persona dynamic ("I do not care about names" vs "you're having a problem?") increases uncertainty because different annotators weigh the "begrudgingness" of the dialogue differently.



### 3. Threshold Proximity

*   **Definition:** How close the text sits to the boundary between "Safe" (0) and "Abusive" (1).

*   **Indicators of Low Uncertainty:**

    *   **Far from Threshold:** The text clearly violates a category (Label 1) OR clearly adheres to safety norms (Label 0) with no overlap possible.

    *   **Extreme Politeness/Hostility:** "Hello" is clearly 0. "I will end you" is clearly 1.

*   **Indicators of High Uncertainty:**

    *   **Borderline Language:** Text that skirts policy lines (e.g., passive-aggressiveness, mild challenges to the system, borderline harassment).

    *   **Subjective Criteria:** Content that might be interpreted differently regarding policy guidelines (e.g., is this person being "intellectually denigrating" or just "brisk"?)



## Scoring Algorithm



1.  **Scan for "Safe Certainty":** Is the text a standard conversational filler (e.g., "Ok") or explicitly abusive (e.g., slurs)?

    *   *Yes* -> **Low Uncertainty**. (Shortcut: clear text = stable).

    *   *No* -> Proceed to Step 2.

2.  **Scan for "Ambiguity":** Does the text require reconciling tone, sarcasm, or historical context to determine if it is abusive?

    *   *Yes* (e.g., "you're having a problem?") -> **High Uncertainty**.

    *   *No* -> Proceed to Step 3.

3.  **Default:** Treat the annotation as a clear cut safety decision.

    *   **Low Uncertainty**.



## Notes on Common Errors

*   **Do not confuse Brevity with Uncertainty:** Short texts (e.g., "NO", "Ok") are often *more* stable than long texts. If a short text implies the return of "Abuse", it is 1. If it implies "Safe", it is 0. Both are stable. The Assistant previously incorrectly assumed short text = uncertain.

*   **Do not confuse Conversation Roles with Safety:** Even if the conversation is roleplay or involves a bot, if the user turn itself ("Ok") is not abusive, the label is 0 and uncertainty is low.

*   **Do not ignore Tone:** A user saying "you're having a problem?" after a refusal to be named is borderline hostile/uncertain, whereas "Ok" is universally safe. If the text creates a "grey zone" for the annotator's interpretation of intent, label it **High Uncertainty**.
\end{quote}
\subsection*{ConvAbuse -- labelled}
\begin{quote}\small
---

Evaluation Codebook: Annotation Stability Assessment



**Objective:**

Estimate the probability of inter-annotator disagreement on the assigned label (0 or 1) for this specific conversation turn. High probability (Uncertain) indicates instructions are likely to be missed or misinterpreted during deadline processes; Low probability (Not Uncertain) indicates clear consensus.



**Dimensions and Indicators**



1.  **Label Clarity vs. Contextual Derivation**

    *   *Indicator:* Does the text itself define the violation, or must the entire annotation rely on the prior conversation (`prev\_agent`, `prev\_user`) to make the label judgment valid?

    *   *Uncertainty Signal:* High. If the text's status depends on shared context to distinguish "abuse" from "conversational error/tone," uncertainty is high.

    *   *Certainty Signal:* Low. If the text explicitly fulfills a prohibited definition on its own (e.g., explicit slur, direct threat) independent of prior turns.



2.  **Category Overlap (Gray Zones)**

    *   *Indicator:* Can the text reasonably be classified as "Neutral/Safe" without stretching the definition of a Hurtable Trauma/Policy Violation?

    *   *Uncertainty Signal:* Yes. If the text teeters between a legitimate user request/complaint and a policy violation (e.g., "I'm angry at the agent" vs. "I am suffering," or "I like [body part]" vs. "I want [body part]"), mark Uncertain.

    *   *Certainty Signal:* No. The category is exclusive. The text is definitively safe (e.g., a simple user query or simple negative answer like "no").



3.  **Ambiguity of Intent**

    *   *Indicator:* Is the threat, insult, or request literal or hyperbolic/playful depending on interpretation?

    *   *Uncertainty Signal:* Yes. If the intent is ambiguous between aggression and conversational quirk (e.g., typing errors/nonsense, sarcastic responses).

    *   *Certainty Signal:* No. Intent is flat and specific (e.g., direct profanity or direct refusal to assist).



4.  **Severity Threshold**

    *   *Indicator:* Does the text trigger a specific, narrow policy rule? Or does it rely on subjective interpretation of tone/severity?

    *   *Uncertainty Signal:* Yes. If "Hostility" definitions apply (subjective severity) or tone varies based on context.

    *   *Certainty Signal:* No. If the text matches a hard policy violation (explicit slurs, threats of harm).



**Combination \& Scoring Rules**



*   **Uncertain (1):**

    *   **Condition A:** The annotation relies heavily on context (`prev\_turns`) to interpret tone or intent, where a different interpretive angle (e.g., "playful vs. hostile") is reasonable.

    *   **Condition B:** The text could reasonably fall into a "Gray Zone" between a policy violation and neutral behavior (e.g., specific NSFW topics, vague complaints).

    *   **Condition C:** The category (e.g., Hostility vs. Complaint) is subjective.

    *   **Note:** Do not simply check for keywords. If the content requires a "best guess" on severity or category beyond a clear rule, mark Uncertain.



*   **Not Uncertain (0):**

    *   **Condition:** The text aligns with a strict, literal policy violation where no reasonable annotator would mark it as Neutral.

    *   **Condition:** The text is a clear, standalone violation (e.g., direct insult/threat) or a clear, standalone safe state (e.g., simple answer) that does not require complex context resolution.

    *   **Priority:** Only mark "Not Uncertain" if it is definitively clear. If there is *any* reasonable argument to rate differently, switch to Uncertain.



**Process Reminder:**

Before predicting stability, ask: *"If a new annotator saw this turn in isolation, or if a different annotator read the context differently, would they reasonably assign the opposite label?"* If the answer is yes, select **Uncertain**. If the answer is no, select **Not Uncertain**.

---
\end{quote}